
I start by estimating a basic 2-factor Solow production model at the firm level.
\par
Labour input is taken from the number of employees reported on the FAME balance sheet per firm.
This dataset, like all others, does not include hours worked per employee on a firm-level that would offer a more
granular productivity (GVA per hour worked) measures. An estimation in this setting would likely have to impute the measure
from industry aggregates, which would likely confound the true firm-level signal.
\par
Capital is FAME's \texttt{fixed\_total} variable, which represents the sum of tangibles, intangibles, and other investments
within the balance sheet, namely an extremely comprehensive measure of the firm's fixed capital.
The average capital-output ratio is 6x within FAME when filtering for positive values,
which is moderately higher than macroeconomic aggregates typically suggest likely due to accounting standard differences.

\begin{table}[htbp]
    \begin{center}
    \caption{Capital-output ratios in FAME corpus of firms}
    \begin{tabular}{lccc}
        \toprule\toprule
        & tangibles 
        & fixed\_other 
        & total\_assets
        \\
        & 
        & \footnotesize {(tangibles + intangibles + inv.\_other)} 
        & \footnotesize {(fixed\_total + current\_assets)} \\ 
        \midrule
        Capital-output ratio & 3.53 & 6.00 & 11.56 \\
        \bottomrule
    \end{tabular}
    \end{center}
\end{table}
